\chapter{Addresses, tiles and cache observations}
\MARSChapterHeading{Addresses, tiles and cache observations}
\label{revision-lab:locality}

Use \texttt{matrix-locality.asm} from \texttt{revision/labs/files}. The complete
program implements original project \texttt{011-matrix-transpose}'s two functions.
Read the main chapter's source/output mapping before inspecting performance.
This lab uses ordinary Run and Debug and the native Data Cache Simulator, followed
by an independent clockwise-rotation task.

\section{Predict addresses in a rectangular matrix}

For the default three-by-five source, locate values five and one by both coordinate
and linear index. Derive their destination coordinates and byte displacements.
Then write all five output rows. Explain why the output row width is three even
though the source row width is five.

For the first source row, derive both address sequences and classify the first ten
read/write notices in a cold fully associative eight-line cache with four-word
lines. Account for separate source and destination bases. Compare your prediction
with the main chapter's trace only after identifying which accesses share a line.

Set a breakpoint at \texttt{transpose\_begin}. Its next instruction calls the
selected function; arguments already contain pointers and dimensions. Use the
Data Cache Simulator in MARS Tools and select the active execution. Expand Cache
Settings and verify eight Cache blocks, four Words per block, Fully associative
placement, eight Blocks per set and Least recently used replacement. The source's
cache annotation supplies these defaults. If changing settings manually, apply
them while stopped at entry, then confirm that the cache is cold and counters zero.

Continue to \texttt{transpose\_done}, before the caller's count store. Inspect
\texttt{\$v0}=15 and output's five rows of three words. The complete ordinary
function should have 30 cache accesses and eight misses. Stepping through the
following caller store adds another data operation; do not compare that later
boundary with the function-only measurement.

\section{Inspect the first clipped tile}

Set \texttt{METHOD} to one in a copy while keeping \texttt{ROWS}=3 and
\texttt{COLS}=5. At the first \texttt{tile\_begin}, identify origin $(0,0)$.
Derive exclusive ends three and four. The first tile has twelve cells; there is
no fourth source row to read.

Continue to the next \texttt{tile\_begin}. Its column origin is four and its row
origin remains zero. The preceding tile is now complete. Output indices zero
through eleven have their final transposed values; indices twelve through fourteen
still have their initial contents. In the source, the completed reads are columns
zero through three of each row, leaving values four, nine and fourteen for the
last tile. Do not mistake the tile for a contiguous source prefix.

With a cold fully associative cache at function entry, this boundary has 24 accesses
and seven misses. Derive those seven resident lines: four source lines and three
destination lines. The final tile adds six accesses and one new destination line,
bringing totals to 30 and eight. Use the stopped source boundary and the cache log
to distinguish an unfinished tile from a completed small matrix.

Make a disposable copy with the row end left at four instead of clipped to three.
Locate the first load outside the active source window. A readable allocation byte
can still be outside the matrix passed to this call. Restore the row clip, then
investigate the column clip separately with five columns.

\section{Compare equal-capacity placements}

Use a sixteen-by-sixteen source and start from a fresh execution for each method.
Keep both allocation addresses and all cache settings fixed except the placement
being investigated. Fully associative ordinary traversal gives 512 accesses and
320 misses. The tiled traversal gives the same 512 accesses and 128 misses.
Explain why equal operation counts can have different line reuse.

For direct mapping, select one block per set and keep eight total blocks. For
two-way placement, keep eight total blocks and use two blocks per set. Confirm
the resulting 128-byte capacity in each view. Apply and reset at the entry boundary,
then execute the function again. Applying settings after the function has returned
does not replay its earlier memory operations.

At this shape, direct mapping gives 344 misses for either method, and two-way
placement gives 320 for either. Derive the row-start sets using the actual addresses.
A destination row stride of 64 bytes advances four sixteen-byte lines. Explain
why several live destination lines compete for only two direct-map sets or one
two-way set, despite other cache entries being available. Figure~\ref{fig:revision-locality-cache}
shows the actual full and direct tiled observations at the same return boundary.

\begin{figure}[htbp]
  \centering
  \begin{tikzpicture}[alt={Two actual native Data Cache Simulator observations after the same sixteen by sixteen tiled transpose, before the caller store. Both use eight four-word blocks, LRU, 128 bytes and 512 accesses. The left uses fully associative placement and eight blocks per set: 384 hits, 128 misses, 75 percent hit rate. The right uses direct mapping and one block per set: 168 hits, 344 misses, 33 percent hit rate. Native publication boxes mark placement and counter rows.}]
    \node[anchor=north west,inner sep=0,font=\small\bfseries] (atitle) at (0,0) {1. Fully associative};
    \node[anchor=north west,inner sep=0] (a) at ([yshift=-2mm]atitle.south west) {\MARSRevisionCrop{revision/captures/locality/transpose-full-tiled.png}{0.40\linewidth}{0.638281250}{0.320000000}{0.185937500}{0.257777778}};
    \coordinate (borigin) at ([xshift=12mm]a.north east);
    \node[anchor=north west,inner sep=0,font=\small\bfseries] (btitle) at (borigin |- atitle.north west) {2. Direct mapping};
    \node[anchor=north west,inner sep=0] (b) at (btitle.south west |- a.north west) {\MARSRevisionCrop{revision/captures/locality/transpose-direct-tiled.png}{0.40\linewidth}{0.638281250}{0.320000000}{0.185937500}{0.257777778}};
    \foreach \panel in {a,b} {
        \begin{scope}[shift={(\panel.south west)},x={($(\panel.south east)-(\panel.south west)$)},y={($(\panel.north west)-(\panel.south west)$)}]
          \draw[black,line width=.7pt] (.025,.83) rectangle (.72,.99);
          \draw[black,line width=.7pt] (.025,.08) rectangle (.96,.22);
        \end{scope}
      }
  \end{tikzpicture}
  \caption{Actual light-theme Debug observations of the same tiled program and matrix. Capacity and access count agree; placement changes the lines that can coexist. Boxes are native LaTeX annotations identifying the placement and counters.}
  \label{fig:revision-locality-cache}
\end{figure}


Repeat with eight-by-eight or fifteen-by-thirteen. Compare geometry before drawing
a conclusion about associativity or tile size. Different dimensions change row
alignment, line occupancy, clipping and set competition. No one miss count is a
correctness requirement for every valid implementation.

\section{Reproduce and interpret the measurement}

Run the public tool from the repository root with JDK 25:
\begin{verbatim}
python3 tools/measure_revision_transpose.py \
  --java /path/to/jdk-25/bin/java \
  --cli /path/to/cli.jar --output /new/report/directory
\end{verbatim}
It checks the transpose and records twelve shape/method combinations. Each case
starts a fresh session and cold full/direct/two-way observations at function entry,
then stops before the caller's count store. Retain its JSON receipt, source and CLI
hashes, stdout and access log. Use a new directory for another run.

For sixteen-by-sixteen, compare ordinary 1,643 instructions with tiled 2,457.
Account for tile-end calculations, per-tile pointer formation and extra loop levels.
The cache improvement is real in the specified fully associative model; the extra
instructions are real too. A claim about elapsed time would require an appropriate
timing experiment or a justified hardware model, neither of which these counters
provide.

Distinguish the experiment's boundaries from the tools' behavior. Instruction
fetch and debugger inspection are excluded from these data counts. Auxiliary
stack traffic inside a different implementation would be included. Backstep does
not rewind the cache observations; replaying an instruction adds another forward
access. Restart or reset at the correct stopped entry for a new cold comparison.

Run the original functional/traversal checks as well:
\begin{verbatim}
java -jar /path/to/cli.jar grade \
  packages/011-matrix-transpose/checks/public.kts \
  /path/to/your/matrix-locality.asm
\end{verbatim}
A correct final output does not prove the one-read/one-write or complete-tile
schedule. For a repeated-load fault, find the duplicated address in the log.
For a wrong rectangular stride, derive the first incorrect destination address.
Fix correctness before interpreting a lower miss count that resulted from skipped
work or accesses outside the active matrix.

\section{Implement clockwise rotation independently}

Open \texttt{matrix-rotate.asm} and implement \texttt{rotate\_clockwise}. Arguments
are source pointer, destination pointer, rows and columns in
\texttt{\$a0} through \texttt{\$a3}. Dimensions are zero through sixteen;
word-aligned nonoverlapping windows lie within the named allocations. Return
$RC$ in \texttt{\$v0}. Source $(i,j)$ must reach output $(j,R-1-i)$.

Predict the complete two-by-three example before coding. Trace source top-left
to output top-right and source bottom-left to output top-left. Explain why swapping
coordinates alone gives a transpose rather than the required clockwise turn.
Choose a source-oriented or destination-oriented traversal, or another valid order;
this task does not require a particular tile schedule.

Keep each active source read and destination write an aligned word notice, exactly
once per cell, and read a cell before its corresponding write. Preserve source,
inactive destination words, guards, saved registers, stack pointer and caller-owned
stack bytes. Do not read the output as temporary workspace. Private stack storage
is permitted under the preservation rules.

Test empty shapes before entering a loop. Test both a single row and a single
column, then a rectangle and an interior window whose pointer differs from its
allocation label. Include words with high bits set: the required operation copies
bit patterns, not just small positive numbers. Run:
\begin{verbatim}
java -jar /path/to/cli.jar grade \
  revision/labs/checks/matrix-rotate.kts \
  /path/to/your/matrix-rotate.asm
\end{verbatim}
The unfinished starter stops explicitly. When a case fails, preserve its shape,
window offsets and first incorrect mapped coordinate. Establish the mapping and
bounds before experimenting with cache scheduling. Once correct, compare the
access patterns of two valid traversals under a stated cache policy; keep that
performance investigation separate from the function's acceptance criteria.
